% !TEX root = chapter-standalone.tex

\section{Monoid actions}
\label{sec:actions}



\begin{ctdefinition}[Monoid action]
    \label{def:monoid-cov-action}
    An \maindef{action} of a \SY{monoid} $\monoidA = \tup{\monoidAset, \mtimesof\monoidA, \idmon_\monoidA}$ on a set $\stateset$ is:

    \constit
    \begin{enumerate}
        \item A \SY{monoid homomorphism}
              \begin{equation}
                  \label{eq:act-monoid}
                  \act \colon \monoidA \fto \Endof\stateset.
              \end{equation}
    \end{enumerate}
\end{ctdefinition}

Unpacked, the above definition means that we have a function
\begin{equation}
     \act \colon  \monoidAset \to (\stateset \to \stateset)
\end{equation}
that satisfies the following two conditions.
 
First, for all $\monelA, \monelB \setin \monoidAset$,
\begin{equation}
    \label{eq:act-to-endo-comp-compatibility}
    \act(\monelA \mtimesof\monoidA \monelB) = \act(\monelA) \mtimesof{\Endof\stateset} \act(\monelB), 
\end{equation}
which encodes ``compatibility of $\act$ with the composition operations in the source and the target'';

Second, 
\begin{equation}\label{eq:act-neutral-element-condit}
    \act(\idmon_\monoidA) = \catidat\stateset,
\end{equation}
which encodes ``compatibility of $\act$ with the neutral elements in the source and the target''. 

Both conditions may be unpacked even further, in terms of individual elements of $\stateset$. 
Recall that the composition operation of $\Endof\stateset$ is the composition of functions, $\mapa \mtimesof{\Endof\stateset} \mapb = \mapa \mthen \mapb$ for $\mapa, \mapb \setin \Endof\stateset$, in which the function written first is applied first (\cref{exa:endofunction-monoid}). 
Hence the first condition says
\begin{equation}
    \label{eq:act-to-endo-comp-compatibility-elementwise}
    \forall \monelA, \monelB \setin \monoidAset \colon \forall \statesetel \setin \stateset \colon \quad \act(\monelA \mtimesof\monoidA \monelB)(\statesetel) = \act(\monelB)\bigl(\act(\monelA)(\statesetel)\bigr).
\end{equation}
In words: acting by $\monelA \mtimesof\monoidA \monelB$ is the same as first acting by $\monelA$ and then acting by $\monelB$.
The second condition says
\begin{equation}
    \forall \statesetel \setin \stateset \colon \quad \act(\idmon_\monoidA)(\statesetel) = \statesetel.
\end{equation}

\todotext{Connect to examples now here...}

\begin{example}[Matrices acting on vectors]
    \label{exa:matrix-action}
    Fix a natural number $n$, and let $\reals^{\ntimesn}$ denote the set of $\ntimesn$ real matrices.
    Consider the \SY{monoid} $\monoidA = \tup{\reals^{\ntimesn}, \mtimesof\monoidA, \idmat}$ whose composition operation is defined, for all $\matmor{A}, \matmor{B} \setin \reals^{\ntimesn}$, by
    \begin{equation}
        \label{eq:matrix-action-composition}
        \matmor{A} \mtimesof\monoidA \matmor{B} \definedas \matmor{B} \matmor{A},
    \end{equation}
    where the right-hand side denotes the usual product of matrices, and whose \SY{neutral element} is the identity matrix $\idmat$.
    This is indeed a \SY{monoid}: associativity follows from the associativity of matrix multiplication, since
    \begin{equation}
        (\matmor{A} \mtimesof\monoidA \matmor{B}) \mtimesof\monoidA \matmor{C} = \matmor{C} (\matmor{B} \matmor{A}) = (\matmor{C} \matmor{B}) \matmor{A} = \matmor{A} \mtimesof\monoidA (\matmor{B} \mtimesof\monoidA \matmor{C}),
    \end{equation}
    and the neutrality laws hold because $\matmor{A} \idmat = \matmor{A} = \idmat \matmor{A}$.
    (In other words, $\monoidA$ is the opposite of the monoid of $\ntimesn$ matrices under the usual matrix multiplication; see \cref{exa:opposite-matrices}.)

    We read $\matmor{A} \mtimesof\monoidA \matmor{B}$ as ``first $\matmor{A}$, then $\matmor{B}$''.
    This matches what matrices do to vectors: to transform a vector $\vectob{v} \setin \reals^n$ first by $\matmor{A}$ and then by $\matmor{B}$, we compute $\matmor{B}(\matmor{A}\vectob{v}) = (\matmor{B}\matmor{A})\vectob{v}$.

    We define an action of $\monoidA$ on the set $\reals^n$ of real $n$-dimensional (column) vectors by
    \begin{equation}
        \label{eq:matrix-action}
        \act \colon \monoidA \fto \Endof{\reals^n},
    \end{equation}
    where, for each $\matmor{A} \setin \reals^{\ntimesn}$,
    \begin{equation}
        \defmapperiod{\act(\matmor{A})}{\reals^n}{\to}{\reals^n}{\vectob{v}}{\matmor{A}\vectob{v}}
    \end{equation}
    Let us check that $\act$ is a \SY{monoid morphism}.
    For all $\matmor{A}, \matmor{B} \setin \reals^{\ntimesn}$ and all $\vectob{v} \setin \reals^n$ we have
    \begin{equation}
        \act(\matmor{A} \mtimesof\monoidA \matmor{B})(\vectob{v})
        = (\matmor{B}\matmor{A})\vectob{v}
        = \matmor{B}(\matmor{A}\vectob{v})
        = \act(\matmor{B})\bigl(\act(\matmor{A})(\vectob{v})\bigr)
        = \bigl(\act(\matmor{A}) \mtimesof{\Endof{\reals^n}} \act(\matmor{B})\bigr)(\vectob{v}),
    \end{equation}
    which is condition \cref{eq:act-to-endo-comp-compatibility-elementwise}, and
    \begin{equation}
        \act(\idmat)(\vectob{v}) = \idmat \vectob{v} = \vectob{v} = \catidat{\reals^n}(\vectob{v}),
    \end{equation}
    which is condition \cref{eq:act-neutral-element-condit}.

    Note that the reversed order in \cref{eq:matrix-action-composition} is essential.
    If we had instead equipped $\reals^{\ntimesn}$ with the usual matrix multiplication, $\matmor{A} \cdot \matmor{B} = \matmor{A}\matmor{B}$, then the same assignment $\matmor{A} \mapsto (\vectob{v} \mapsto \matmor{A}\vectob{v})$ would satisfy
    \begin{equation}
        \act(\matmor{A}\matmor{B}) = \act(\matmor{B}) \mtimesof{\Endof{\reals^n}} \act(\matmor{A}),
    \end{equation}
    with the order reversed, and so, for $n \geq 2$, it would not be a \SY{monoid morphism} into $\Endof{\reals^n}$.
\end{example}

\begin{example}[Scalar multiplication]
Let $\monoidA = \tup{\reals, \cdot, 1}$ be the monoid of real numbers, with multiplication as composition operation and the number $1$ as neutral element. Consider the set $\stateset = \reals^3$ of real 3-dimensional vectors. An action 
\begin{equation}
\act \colon \reals \fto \Endof{\reals^3}
\end{equation}
is given by scalar multiplication: for each $\lambda \in \reals$, 
\begin{equation}
\defmapperiod{\act(\lambda)}{\reals^3}{\to}{\reals^3}{\begin{pmatrix} x_1 \\ x_2 \\ x_3 \end{pmatrix}}{\begin{pmatrix} x_1 \cdot \lambda \\ x_2 \cdot \lambda \\ x_3 \cdot \lambda \end{pmatrix}}
\end{equation}
\end{example}

\begin{example}[A memory bit: the flip-flop monoid]
    \label{exa:flip-flop-action}
    A memory cell that stores one bit (for example, a latch in a digital circuit) can receive two commands: $\text{set}$, which stores the value $1$, and $\text{reset}$, which stores the value $0$.
    Together with the ``empty command'' $\idmon$, which does nothing, these commands form the \emph{flip-flop monoid}
    \begin{equation}
        \mathbf{F} = \tup{\makeset{\idmon, \text{set}, \text{reset}}, \mtimesof{\mathbf{F}}, \idmon},
    \end{equation}
    in which $\monelA \mtimesof{\mathbf{F}} \monelB$ means ``first $\monelA$, then $\monelB$''.
    Since the cell only remembers the last value written to it, its composition operation is ``the last command wins'':
    \begin{equation}
        \monelA \mtimesof{\mathbf{F}} \monelB \definedas
        \begin{cases}
            \monelB & \text{if } \monelB \neq \idmon, \\
            \monelA & \text{if } \monelB = \idmon.
        \end{cases}
    \end{equation}
    For instance, $\text{set} \mtimesof{\mathbf{F}} \text{reset} = \text{reset}$, whereas $\text{reset} \mtimesof{\mathbf{F}} \text{set} = \text{set}$; so $\mathbf{F}$ is not commutative.
    It is a monoid: $\idmon$ is neutral by definition, and both $(\monelA \mtimesof{\mathbf{F}} \monelB) \mtimesof{\mathbf{F}} \stylemorph{p}$ and $\monelA \mtimesof{\mathbf{F}} (\monelB \mtimesof{\mathbf{F}} \stylemorph{p})$ are equal to the last element among $\monelA, \monelB, \stylemorph{p}$ that is different from $\idmon$ (or to $\idmon$, if there is none).
    Note also that $\text{set} \mtimesof{\mathbf{F}} \text{set} = \text{set}$: writing the same value twice is the same as writing it once.

    The monoid $\mathbf{F}$ acts on the set $\stateset = \makeset{0, 1}$ of possible contents of the cell by
    \begin{equation}
        \act(\idmon) = \catidat{\stateset}, \qquad
        \act(\text{set}) = (\statesetel \mapsto 1), \qquad
        \act(\text{reset}) = (\statesetel \mapsto 0).
    \end{equation}
    To check that $\act(\monelA \mtimesof{\mathbf{F}} \monelB) = \act(\monelA) \mthen \act(\monelB)$, note that if $\monelB \neq \idmon$, then $\act(\monelB)$ is a constant function, so that $\act(\monelA) \mthen \act(\monelB) = \act(\monelB)$, and if $\monelB = \idmon$, then $\act(\monelA) \mthen \act(\monelB) = \act(\monelA)$.
    In both cases, this is $\act(\monelA \mtimesof{\mathbf{F}} \monelB)$.

    More generally, a storage cell for values from any set $\stylesets{V}$ (a ``register'') is described by the monoid whose elements are $\idmon$ and one command ``write $v$'' for each $v \setin \stylesets{V}$, again with ``the last command wins'', acting on $\stylesets{V}$.
    The flip-flop monoid will reappear later: in the change detector of \cref{exa:mealy-change-detector}, the inputs $0$ and $1$ act on the states exactly as $\text{reset}$ and $\text{set}$.
    Together with groups, it is one of the basic building blocks of the theory of finite automata (Krohn--Rhodes theory).
\end{example}

\begin{example}[Charging a battery: a saturating counter]
    \label{exa:saturating-counter}
    The charge of a battery is often displayed in a few discrete levels, say $\stateset = \makeset{0, 1, \dots, n}$, where $0$ means empty and $n$ means full (for instance, $n = 4$ bars on a phone display).
    Charging the battery for $a$ units of time raises its level by $a$, except that the level cannot exceed $n$: the new level is $\min(\statesetel + a, n)$.
    In particular, charging for $n$ units or more always results in a full battery, whatever the initial level.
    So the possible ``charging commands'', up to their effect, are $0, 1, \dots, n$, and composing two of them is addition, capped at $n$.
    This is the \emph{saturating counter} monoid
    \begin{equation}
        \mathbf{C}_n = \tup{\makeset{0, 1, \dots, n}, \oplus, 0}, \qquad a \oplus b \definedas \min(a + b, n).
    \end{equation}
    It is indeed a monoid: $0$ is neutral, and both $(a \oplus b) \oplus c$ and $a \oplus (b \oplus c)$ are equal to $\min(a + b + c, n)$.
    It acts on the set $\stateset$ of charge levels by charging:
    \begin{equation}
        \act(a, \statesetel) \definedas \min(\statesetel + a, n).
    \end{equation}
    The condition for the neutral element holds since $\min(\statesetel, n) = \statesetel$, and compatibility with composition holds since
    \begin{equation}
        \begin{aligned}
            \act(b, \act(a, \statesetel))
            & = \min\bigl(\min(\statesetel + a, n) + b, n\bigr) \\
            & = \min(\statesetel + a + b, n) = \act(a \oplus b, \statesetel).
        \end{aligned}
    \end{equation}
    For example, for $n = 4$, charging a battery at level $1$ first by $2$ units and then by $3$ units gives $\act(3, \act(2, 1)) = \act(3, 3) = 4$, the same as charging it by $2 \oplus 3 = 4$ units at once.

    In the same way, \emph{discharging}, $\tup{a, \statesetel} \mapsto \max(\statesetel - a, 0)$, is another action of $\mathbf{C}_n$ on $\stateset$.
    Charging and discharging do not commute, however: for $n \geq 2$, charging a full battery by $2$ and then discharging it by $2$ leaves it at level $n - 2$, whereas discharging first and then charging leaves it full.
    So, to describe both kinds of commands at once, we would need a monoid which is not commutative.
\end{example}

\begin{example}[Relations acting on sets of possible states]
    \label{exa:relations-acting-on-subsets}
    Many systems are \emph{nondeterministic}: from a given state, a step may lead to one of several possible states.
    For example, a robot moving along a corridor with positions $\stateset = \makeset{1, 2, 3}$ may receive the command ``forward'', but its wheels sometimes slip, so that it moves forward by one position or stays where it is (and it stays at the end of the corridor, position $3$).
    Such a step is described by a relation $\relA \colon \stateset \mto \stateset$ (\cref{def:binary-relation}), where $\inrel{\statesetel}{\relA}{\statesetel'}$ means that $\statesetel'$ is a possible next state when the current state is $\statesetel$; for the slipping robot,
    \begin{equation}
        \relA_{\text{fwd}} = \makeset{\tup{1, 1}, \tup{1, 2}, \tup{2, 2}, \tup{2, 3}, \tup{3, 3}}.
    \end{equation}
    The relations $\stateset \mto \stateset$, with the composition of relations (\cref{def:composition-of-relations}) and the identity relation (\cref{def:identity-relation}), form a monoid.
    If $\stateset$ is finite, this monoid is finite as well: for $\stateset$ with $n$ elements, there are $2^{n^2}$ relations.

    If we do not know the current state exactly, but only a set $\stylesets{P} \subseteq \stateset$ of possible states, then after a step described by $\relA$ the set of possible states is
    \begin{equation}
        \act(\relA)(\stylesets{P}) \definedas \makeset{\statesetel' \setin \stateset \mid \exists \statesetel \setin \stylesets{P} \colon \inrel{\statesetel}{\relA}{\statesetel'}}.
    \end{equation}
    This defines an action of the monoid of relations on the set $\powerset(\stateset)$ of subsets of $\stateset$.
    Indeed, the identity relation does not change $\stylesets{P}$, and $\act(\relA \mthen \relB)(\stylesets{P}) = \act(\relB)(\act(\relA)(\stylesets{P}))$, since a state is reachable from $\stylesets{P}$ by a step of $\relA \mthen \relB$ exactly if it is reachable by a step of $\relA$ followed by a step of $\relB$.
    For the slipping robot starting at position $1$, the sets of possible positions after zero, one and two commands ``forward'' are
    \begin{equation}
        \makeset{1}, \qquad \makeset{1, 2}, \qquad \makeset{1, 2, 3}.
    \end{equation}

    In terms of the matrices of \cref{ex:ComposingFiniteRelations}, a relation on $\stateset = \makeset{1, \dots, n}$ is an $\ntimesn$ matrix $\phi_{\relA}$ with entries in $\makeset{0, 1}$, and a subset $\stylesets{P}$ is a row vector $\stylemaps{p}$ with entries in $\makeset{0, 1}$.
    Composition of relations is the ``Boolean'' matrix product, in which sums are replaced by maxima, and the action is the Boolean product $\stylemaps{p} \, \phi_{\relA}$ of a row vector with a matrix.
    Since the row vector is multiplied from the left, the order of composition comes out right without reversing it: acting first by $\phi_{\relA}$ and then by $\phi_{\relB}$ gives $(\stylemaps{p} \, \phi_{\relA}) \, \phi_{\relB} = \stylemaps{p} \, (\phi_{\relA} \, \phi_{\relB})$ (compare \cref{exa:matrix-action}).
    For the slipping robot,
    \begin{equation}
        \phi_{\relA_{\text{fwd}}} = \begin{bmatrix} 1 & 1 & 0 \\ 0 & 1 & 1 \\ 0 & 0 & 1 \end{bmatrix},
        \qquad
        \begin{bmatrix} 1 & 0 & 0 \end{bmatrix} \phi_{\relA_{\text{fwd}}} = \begin{bmatrix} 1 & 1 & 0 \end{bmatrix}.
    \end{equation}
\end{example}

\devel{
\begin{example}[robot on a grid]
\todotext{connect the intro to the definition}
\end{example}



\todotext{Make a remark/example somewhere that any monoid acts on itself...}

}





\subsection{Uncurrying actions}

Recall that (right) currying is the operation of applying the canonical isomorphism
\begin{equation}
    \curry \colon \setC^{\setA \cartprod \setB} \to \pars{\setC^\setB}^\setA
\end{equation}
which transforms any function of the type
\begin{equation}
    \setA \cartprod \setB \to \setC
\end{equation}
into one of the type
\begin{equation}
   \setA \to (\setB \to \setC).
\end{equation}

A monoid action involves a function of the type
\begin{equation}
     \act \colon  \monoidAset \to (\stateset \to \stateset).
\end{equation}
If we apply the inverse $\curry^{-1}$ (``uncurrying'') to this function, 
we obtain a function of the type
\begin{equation}\label{eq:uncurried-action}
 \makecartprod{\monoidAset, \stateset} \sto \stateset.
 \end{equation}
Throughout mathematics and applications, actions are often presented in this uncurried form. 
Since we can tell the difference between $\act$ and its uncurried version based on their respective domains and codomains, we will use the same name, ``$\act$'', for both versions.


Now, so far, we have just used $\curry^{-1}$ to produce the function \cref{eq:uncurried-action} from a monoid action. However, $\act \colon  \monoidAset \to (\stateset \to \stateset)$ was not just any function, it was a monoid homomorphism $\act \colon \monoidA \fto \Endof\stateset$. What does that condition translate to in the uncurried version $\act \colon  \makecartprod{\monoidAset, \stateset} \sto \stateset$? The answer is summarized in the following definition.

\begin{ctdefinition}[Monoid action, uncurried]
\label{def:uncurried-action}
   
    An \emph{uncurried action} of a \SY{monoid}~$\monoidA = \tup{\monoidAset, \mtimesof\monoidA, \idmon_\monoidA}$ on a set~\stateset is:
    
    \constit
    \begin{enumerate}
        \item a function
              \begin{equation}
                  \label{eq:ract-monoid}
                  \act \colon \makecartprod{\monoidAset, \stateset} \sto \stateset.
              \end{equation}
    \end{enumerate}
    \condit
    \begin{enumerate}
        \item for all~$\statesetel \setin \stateset$ and $\monelA,\monelB\setin\monoidAset$,
              \begin{equation}
                  \label{eq:monoiad-ract-1cond}
                  \act(\monelA \mtimesof\monoidA \monelB, \statesetel) = \act(\monelB, \act(\monelA, \statesetel));
              \end{equation}
        \item for all~$\statesetel \setin \stateset$,
              \begin{equation}
                  \label{eq:monoiad-ract-2cond}
                  \act(\idmon_\monoidA, \statesetel) = \statesetel.
              \end{equation}
    \end{enumerate}
\end{ctdefinition}

\begin{gradedexercise}[\exname{UncurriedAction}]
    \label{ex:UncurriedAction}
Throughout this exercise, fix a monoid $\monoidA = \tup{\monoidAset, \mtimesof\monoidA, \idmon_\monoidA}$ and~a set $\stateset$. 
\begin{enumerate}
\item Given a function
    \begin{equation}
    \act \colon \makecartprod{\monoidAset, \stateset} \sto \stateset
    \end{equation}
    that satisfies \cref{def:uncurried-action}, prove that the curried function
    \begin{equation}
    \curry(\act) \colon  \monoidAset \to (\stateset \to \stateset)
    \end{equation}
    satisfies conditions \cref{eq:act-to-endo-comp-compatibility} and \cref{eq:act-neutral-element-condit} and is hence a monoid action. 
    \item Given a monoid action 
    \begin{equation}
    \act \colon  \monoidAset \to (\stateset \to \stateset)
    \end{equation}
    prove that its uncurried version 
    \begin{equation}
    \curry^{-1}(\act) \colon  \makecartprod{\monoidAset, \stateset} \sto \stateset
    \end{equation}
    satisfies \cref{def:uncurried-action}. 
\end{enumerate}
\end{gradedexercise}
\solutionof{UncurriedAction}

\devel{

\todotext{Rewrite this part to discuss the possibility of writing uncurried action with $\monoidAset$ on the right...}

\begin{remark}
    Left actions are often denoted using infix notation of the kind
    \begin{equation}
        \makecartprod{\monoidAset, \stateset} \sto \stateset, \ \tup{\monela, \statesetel} \mapsto \monela \cdot \statesetel
    \end{equation}
    or even simply
    \begin{equation}
        \makecartprod{\monoidAset, \stateset} \sto \stateset, \ \tup{\monela, \statesetel} \mapsto \monela \statesetel
    \end{equation}
    (in regard to the latter, think of the notation often used for multiplication of matrices with vectors).
    If we write our notation ``$\monela \mtimesof\monoidA \monelb$'' as ``$\monela \star \monelb$'' (or even simply ``$\monela \monelb$''), then REF reads as
    \begin{equation}
        \monela \cdot (\monelb \cdot \statesetel) = (\monela \star \monelb) \cdot \statesetel,
    \end{equation}
    or simply
    \begin{equation}
        \monela (\monelb \statesetel) = (\monela  \monelb)  \statesetel.
    \end{equation}
    In other words, we can view REF as a kind of ``associative law'', but where two of the three elements come from $\monoidAset$, while one of the elements comes from $\stateset$.
\end{remark}
}






    \subsection{Contravariant actions}
    
    
    \begin{ctdefinition}[Contravariant monoid action]
        \label{def:monoid-cont-action}
        A \maindef{contravariant action} of a \SY{monoid} $\monoidA = \tup{\monoidAset, \mtimesof\monoidA, \idmon_\monoidA}$ on a set $\stateset$ is:

        \constit
        \begin{enumerate}
            \item A \SY{monoid morphism}
                  \begin{equation}
                      \label{eq:cont-act-monoid}
                      \act \colon \monoidA\op \fto \Endof\stateset.
                  \end{equation}
        \end{enumerate}
    \end{ctdefinition}

  

    % Diagram for the two conditions of def:monoid-uncurried-cont-action (corrected 2026-10-14), commented out for now:
    % \begin{marginfigure}
    %     \centering
    %     \includesag{semigroup-contrav-action}
    %     \caption{The conditions for an uncurried contravariant monoid action.}
    %     \label{fig:semigroup-contrav-action}
    % \end{marginfigure}
    
    
    \begin{ctdefinition}[Contravariant monoid action, uncurried]
        \label{def:monoid-uncurried-cont-action}
        An \emph{uncurried contravariant action} of a \SY{monoid} $\monoidA = \tup{\monoidAset, \mtimesof\monoidA, \idmon_\monoidA}$ on a set $\stateset$ is:

        \constit
        \begin{enumerate}
            \item a function
                  \begin{equation}
                      \label{eq:uncurried-cont-action}
                      \act \colon \makecartprod{\monoidAset, \stateset} \sto \stateset.
                  \end{equation}
        \end{enumerate}
        \condit
        \begin{enumerate}
            \item for all $\statesetel \setin \stateset$ and $\monelA, \monelB \setin \monoidAset$,
                  \begin{equation}
                      \label{eq:uncurried-cont-action-1cond}
                      \act(\monelA \mtimesof\monoidA \monelB, \statesetel) = \act(\monelA, \act(\monelB, \statesetel));
                  \end{equation}
            \item for all $\statesetel \setin \stateset$,
                  \begin{equation}
                      \label{eq:uncurried-cont-action-2cond}
                      \act(\idmon_\monoidA, \statesetel) = \statesetel.
                  \end{equation}
        \end{enumerate}
    \end{ctdefinition}

    \begin{remark}
        Just as in the covariant case (\cref{def:monoid-cov-action}, \cref{def:uncurried-action}, and \cref{ex:UncurriedAction}), one can prove that \cref{def:monoid-cont-action} and \cref{def:monoid-uncurried-cont-action} are equivalent: currying and uncurrying give a one-to-one correspondence between contravariant actions $\act \colon \monoidA\op \fto \Endof\stateset$ and uncurried contravariant actions $\act \colon \makecartprod{\monoidAset, \stateset} \sto \stateset$.
        The proof is the same as for \cref{ex:UncurriedAction}, with $\monoidA$ replaced by $\monoidA\op$: since $\monelA \mtimesof{\monoidA\op} \monelB = \monelB \mtimesof\monoidA \monelA$, the order of the two elements on the left-hand side of \cref{eq:monoiad-ract-1cond} is reversed, which gives \cref{eq:uncurried-cont-action-1cond}.
    \end{remark}



    In contrast to \cref{eq:monoiad-ract-1cond}, \cref{eq:uncurried-cont-action-1cond} says that, if we first act by $\monelB$ to obtain $\act(\monelB, \statesetel)$, and then act by $\monelA$, the result is the same as acting by $\monelA \mtimesof\monoidA \monelB$ (note the inverted order).

    % And what about contravariant actions?
    A \SY{contravariant semigroup action} can be defined as a covariant action of \emph{the opposite} \SY{semigroup}.

   




